% Pure Math Survey 1.8.1; standard edition.
% Advisory reading-review threshold: 8 total pages; REVIEW_NEEDED is non-blocking.
% Only an explicit user max_pages sets a page-count limit; preserve necessary mathematics.
% Reader task: Separate established ranges, counterexamples, method limitations, and authenticated remaining targets.
% Develop the same mathematical core at its bottlenecks; do not add a generic depth appendix.
% Reconstruct definitions and natural variation domains before selecting answers.
% Distinguish intrinsic data, representatives, unknowns, and gauges.
% Put the question's ordered quantifiers before its principal answers.
% Discover source-based frontier candidates beyond the initial bibliography.
% Map each substantive body task: local definitions -> usable result -> proof/import -> use.
% These slots illustrate roles, not a number, proportion, or section requirement.
% Delete or move an inapplicable slot; never invent a result to fill it.
% Keep important hypotheses and conclusion formulas inside statements.
% Nearby long definitions and standard terms may be referenced; avoid needless repetition.
% A formula reference may aid navigation, not replace a material condition or conclusion.
% A separately placed proof needs forward/backward locators, not duplicated results.
% A printed proof sketch must not be recorded as a full proof.
\documentclass[12pt,reqno]{amsart}
\usepackage{math-review}
\title{[INSERT: Topic]: Boundaries and Problems}
\author{}
\date{[INSERT: Checked literature cutoff]}
\newcommand{\PrincipalStatement}{%
\placeholder{Specify the objects, quantified domains, and necessary hypotheses.}
% Use an itemized list only for genuinely independent substantial assumptions.
\placeholder{Write the decisive criterion and conclusion formulas here, with the exact source or status.}
}
% For a useful distant recall, reuse this body and the original number:
% \begin{theoremrecall}{thm:principal}\PrincipalStatement\end{theoremrecall}
% Do not print an immediate duplicate or an automatic second principal theorem.
\begin{document}
% Before drafting: complete problem_formulation and survey_selection in the
% existing reader evidence. Review every principal answer in both directions.
\begin{abstract}
\placeholder{State the selected question and actual answer or boundary.}
\end{abstract}
\maketitle

\section{Introduction}\label{sec:introduction}
\begin{definition}\label{def:question-data}
\placeholder{Define indispensable objects and notation; distinguish fixed data, admissible variations, and unknowns.}
\end{definition}
\placeholder{Pose the selected question with its ordered quantifiers. Justify a substantive restriction or relationship; omit trivial transitions.}
\begin{theorem}\label{thm:principal}
\PrincipalStatement
\end{theorem}
% Change the environment to match its logical role. Add other principal answers
% only when selected by the question. Put secondary consequences after the result.
\placeholder{Locate the proof or give the controlling primary source and theorem locator. A citation does not replace the statement above.}

% Activate a Problem only after authenticating its formulation and status.
% \begin{problem}\label{prob:target}
% [State objects, hypotheses, quantifiers, and the exact target conclusion.]
% \end{problem}
% Give its primary source, checked status, known range, and remaining range.
% Delete the discussion and references with an unused Problem slot.

\section{[INSERT: The established range]}\label{sec:range}
\begin{definition}\label{def:range}
\placeholder{Define the parameter or regularity boundary and the solution concept used to describe it.}
\end{definition}
\begin{theorem}\label{thm:range}
\placeholder{State the selected established boundary result with all material range, strictness, and regularity conditions and its actual conclusion.}
\end{theorem}
\paragraph{Source and use.}
\placeholder{Give its controlling source and theorem locator. Explain exactly which portion of the organizing question it resolves.}
\placeholder{Develop why the same decisive hypothesis gives this range. Verify an application only when it tests that hypothesis, keeping any imported proof distinct from a local calculation.}

\section{[INSERT: The obstruction or failed extension]}\label{sec:obstruction}
\begin{proposition}\label{prop:obstruction}
\placeholder{Specify a counterexample or obstruction, with its hypotheses and exact negative conclusion. Do not equate failure of a method with nonexistence.}
\end{proposition}
\begin{proof}\label{prf:obstruction}
\placeholder{Work through the same obstruction or counterexample, including the admissibility check and the failed inference. Explain why it does not contradict the established positive result.}
\end{proof}
% Delete the slot if no independent negative result is selected. A context-only
% status discussion needs no artificial lemma, theorem, or counterexample.

\section{[INSERT: The precisely remaining range]}\label{sec:remaining}
\placeholder{Contrast the proved range with the authenticated target using explicit quantifiers. State source-qualified status; an unsuccessful search does not prove openness.}
% Omit this section when the selected scope is settled, or explain that settled
% scope directly. Do not invent a Problem, application, or unresolved claim.

\templatereferences
\end{document}
